%% ========================================================================
%% Prakata
%% ========================================================================
%% STATUS: SCAFFOLD ONLY - paragraf di bawah adalah placeholder generik;
%% lihat docs/book-plan.md untuk positioning final buku ini.
%% ========================================================================

\chapter*{Prakata}
\addcontentsline{toc}{chapter}{Prakata}

Puji syukur kehadirat Tuhan Yang Maha Esa atas segala rahmat dan karunia-Nya
sehingga buku \textit{Laravel Praktis: Membangun Aplikasi Kasir dari Nol}
ini dapat diselesaikan.

Pengembangan aplikasi web adalah keterampilan inti bagi siapa pun yang ingin
membangun perangkat lunak yang dipakai langsung oleh banyak orang setiap hari.
Buku ini dirancang untuk membawa pembaca dari langkah paling awal, yaitu
memahami arsitektur web modern dan pola MVC, hingga membangun sebuah aplikasi
utuh yang tervalidasi, teruji, dan mengekspos REST API siap produksi.

\section*{Tujuan Buku}

Buku ini bertujuan untuk:

\begin{itemize}
    \item Memberikan pemahaman mendalam tentang arsitektur web modern dan
          kerangka kerja Laravel untuk membangun aplikasi berbasis data
    \item Mengajarkan desain basis data, migrasi, dan Eloquent ORM sebagai
          fondasi pengelolaan data aplikasi
    \item Membawa pembaca dari routing dasar hingga aplikasi utuh yang
          menerapkan autentikasi, otorisasi, dan pengolahan data berskala
    \item Menjembatani teori dengan praktik melalui satu studi kasus yang
          dibangun bertahap dari fondasi hingga REST API siap produksi
\end{itemize}

\section*{Struktur Buku}

Buku ini disusun ke dalam empat bagian yang mengalir secara bertahap:

\begin{description}
    \item[Bagian I: Fondasi Laravel dan Arsitektur Web]: Memperkenalkan
          ekosistem Laravel, protokol HTTP, pola MVC, serta antarmuka dengan
          Blade, Tailwind, dan Alpine.
    \item[Bagian II: Data dan Eloquent]: Mengajarkan desain basis data,
          migrasi, seeding, Eloquent ORM dan relasi, hingga validasi dan
          keamanan input.
    \item[Bagian III: Keamanan, Laporan, dan API]: Mencakup autentikasi,
          otorisasi berbasis peran, pengolahan data (impor, ekspor, antrean),
          hingga perancangan dan pembangunan REST API.
    \item[Bagian IV: Menuju Aplikasi Siap Produksi]: Membahas pengujian
          aplikasi, optimasi dan deployment, serta studi kasus menyeluruh
          yang menyatukan seluruh komponen menjadi satu aplikasi utuh.
\end{description}

Studi kasus \textbf{Simple POS}, sebuah aplikasi kasir sederhana dengan
pengelolaan produk, transaksi, laporan, dan API untuk kasir mobile, menemani
pembahasan di setiap bab agar setiap konsep selalu terhubung ke contoh yang
nyata.

\vspace{1cm}

\noindent
Dian Hanifudin Subhi dan Ade Ismail

\cleardoublepage
